\documentclass[10pt,a4paper]{book}
\usepackage[utf8]{inputenc}
\usepackage{amsmath}
\usepackage{graphicx}
\usepackage[top=1.55cm, bottom=2.29cm, left=1.6cm, right=1.47cm]{geometry}

\pagestyle{empty}
\usepackage{tikz}

%algunos símbolos matemáticos y paquete para usar subimágenes
\usepackage{amsmath}
\usepackage{amsfonts}
\usepackage{amssymb}
\usepackage{graphicx}
\usepackage{subfigure}
\usepackage{listings}
\usepackage{appendix}

%Hyperrefs without borders
\usepackage[hidelinks]{hyperref}


%parametros del documento (sus propiedades)
\hypersetup{
    pdftitle={Irene Sanz Nieto - In-hand object recognition and on-line learning using 2D and 3D information - 2014},
    pdfsubject={Bachelor's Thesis - 2014},
    pdfauthor={Irene Sanz Nieto},
    pdfkeywords={palabraclave1} {palabraclave2} {palabraclave3},
    colorlinks,
    citecolor=black,
    filecolor=black,
    linkcolor=black,
    urlcolor=black,
}


\begin{document}

%elementos antes del trabajo en sí se meten dentro de frontmatter
\frontmatter

%cada incluye referencia a un archivo de tipo .tex
\begin{titlepage}
	\begin{center}

		%forma de introducir imágenes. el \\[0.5 cm] de final de línea introduce un salto de ese tamaño.
		%width=1\textwidth indica el tamaño de la imágen (valores entre 0-1). 
		 \includegraphics[size=0.8]{img/uc3m.eps}  \\[0.5 cm]

		\large \textsc{Systems Engineering and Automation Department} \\ [1 cm]

		\large Bachelor's Thesis\\[1 cm]
		\Huge \textbr{OCULAR} \\[1 cm]
		\huge \textsc{In-hand on-line object recognition and learning using 2D and 3D information}\\[8 cm]


		%flushleft alinea a la izquierda el texto
		\begin{flushright} \Large
			\emph{Author:} Irene Sanz Nieto\\[0.5 cm]
			\emph{Supervisor:} Víctor González Pacheco \\
		\end{flushright}

		%rellena de blanco el resto de la página para escribir abajo del todo
		\vfill

		% Bottom of the page
		{\large Leganés, June 2014}

	\end{center}
\end{titlepage}


%\include{licencia}

%\include{evaluacion}

%\include{agradecimientos}

%\include{resumen}


%%%%%% TABLE OF CONTENTS %%%%%%
\begin{huge}
{title & whatever}
\end{huge}



%%%%%% TABLE OF CONTENTS %%%%%%
\newpage
\tableofcontents

%%%%%% TABLE OF FIGURES %%%%%%
\newpage
\listoffigures
\newpage

%%%%%% INTRO %%%%%%
%\addcontentsline{toc}{part}{Introduction}
\part{Introduction}


%	
\section{Motivation}

Technology has evolved enormously in the past years. 
% In particular, robotics has changed and moved from controlled spaces such as factories to human inhabited spaces. 
% In section \ref{context} the reasons behind this shift in the robot's location and function were presented. 
% It was also estated that the importance of the perception systems augmented with the inclusion of robots in human-inhabited areas. 
% % This fact increases the importance of the perception systems being integrated. 
% The correct recognition of objects, persons, areas and situations is key for assistive and social tasks. 
Nowadays most of the robots being developed are only able to recognize a small part of their environment. 
% Most computer vision systems are currently used to navigate between points or to locate certain objects and grasp them. 
The decision algorithm of the robots is normally based on the instructions received from the user. 
Those commands are usually given by voice or inputting the desired task on a certain User Interface (UI).
\\

% For the robots to act as aiding personnel to take care of elder or sick people, the evolution of this decision algorithm is crucial. 
The algorithm decision of the robots must be upgraded if they are to perform tasks in a human inhabited area. %act as aiding personnel to take care of elder or sick people. 
They must be able to respond to commands, but also be aware of their environment and respond autonomously to it. 
% As an example, robots must be able to recognize the danger involved in certain objects or situations. 
% Also, they should perceive the environment and the interactions between humans to discern if, for example, they are arguing or if the conversation is not a beneficial one for the patient. 
% This latter example may be seen in the case of an anorexic person who is talking about food and weight with another person. 
% Then, the robot may intervene changing the subject subtly or trying to end the conversation. 
When this change has occurred, the robots may successfully perform the assistive tasks now reserved only to humans. 
But there is still a long way to research, mainly around the perception of the environment. 
The investment needed for its development is now being held due to the recession that appeared in the late 2000s decade, whose effects are nowadays still present. 
Nevertheless, this lack of funding might be mitigated using research, open data and open-source code. 
\\

I strongly believe in the ideals presented by the Open Source Initiative. 
This project is intended to be an open source code that can be a building block for other researchers that work on robot perception.
There are various scientists that have already studied the importance of the objects around the human in the action recognition. 
For example, in \cite{Delaitre}, a study of the human-object interaction in still images was performed in order to relate those objects to actions. 
	This relation between object and action may also be seen in \cite{Fathi}, in which wearable cameras are used to retrieve the input. 
	The objects being hand-held are recognized and the action associated to them is learned. 
	In this thesis I present a system that allows to easily learn and recognize hand-held objects in real time. 
	It is intended to be the previous step to the association between the object and the action. 



% The economic situation described above forced many experienced professionals to trying their luck creating start-ups. Many of the ideas of those enterprises are having nowadays a huge impact on the society. 
% \\

% One of these open-source projects that appeared was the low-cost 3D printers. These machines have changed the manner of investigating many fields, since they allow to design different pieces easily and have a 3D reproduction in a few hours. 
% \\

% Initially, they were used in investigation and more specifically in robotics, but now they are used in many different fields. Among those fields, there are medicine, construction or even food making. 
% \\

% In medicine, they have been a revolution since they allow to create customized and precise pieces in very few time. They have been used for prosthesis and implants for persons of various ages, even for babies. In the prosthesis fields in particular, the 3D printing technology is being a complete revolution. Before, the prosthesis were very expensive and permitted only fixed movements and combinations. The adaptations to each individual were made in the final product itself, trying to make it as comfortable as possible for the wearer. Nowadays, the prosthesis are customized for each patient, reducing the inconveniences and increasing their usability. Also, they can be easily and cheaply adapted for children as an example, who are still experimenting many changes in their bodies. They are much cheaper than they were before, and everyone with a 3D printer may construct one. 
% In order to 3D print a piece a file with its description is needed. There are many web-pages that store open-source designs that ranges from decoration models to complex prosthesis. This fact is decisive because there is not needed a huge amount of knowledge or money to improve the life quality of a person using these technologies. 
% \\

% There are numerous open-source projects and developers that put in common their knowledge to improve the technology being used. I have used many of them in the previous years, to learn about 3D printing, robotics or programming among other fields. 
% \\

% It is a fact that acquiring knowledge would be much difficult if the Open Source initiative has not been invented. This impulsed me towards developing something useful and that could be used by other people. The idea of creating a software that could be used in robotics investigation but also help people at the same time. 
% \\[0.5cm]

% Many of the projects are aimed at aiding physically impaired people, creating sternal skeletons and robotic arms that could aid them. But a personal fact led me to realize that visually impaired people were not having as much attention. The applications developed for them are still rough to use and also it is difficult for a grown person to develop his remaining senses to supply the information lost. 
% \\

% Besides, in the robotic field new lines of investigation have appeared. The social robots are now a reality and in the near future we will interact with them everyday. In order to understand the human behavior, the recognition of the objects being handled by them is crucial. 
% \\[0.5cm]


% Computer vision has experience an important improvement in the last years through the upgrade of the hardware and software that compose it. The hardware such as acquisition elements (cameras, depth sensors, etc) and computing elements (PCs or other programmable devices) have experienced a rapid advance in the past years. It allowed to process more data that is now obtained more accurately and with less noise. This increase in the computing power of the equipment created a possibility of introducing more complex libraries and frameworks and even operating systems. 
% \\

% Now, the technology is available to solve the problems presented, the aid of visually impaired people and the introduction of new information in the social robotics field. This is how the idea behind this thesis appeared: the creation of a modular software that implements an in-hand object recognition algorithm. 
 
%	%\addcontentsline{toc}{chapter}{Socio-economical context}
\section{Socio-economical context}

%Socio-economic - relating to both social and economic factors (social groups and the class system for example)
%Context - The circumstances/environment/events surrounding a specific thing. 



%	%\addcontentsline{toc}{chapter}{Future applications}
\chapter{Future applications}
	% future applications



\chapter{Introduction}
This bachelor's thesis consists on a software that implements an in-hand object learning and recognition using 2D and 3D information. 
\\

The code is Open-Source and it was designed to be running inside a robot. The idea is to create a new stand-alone module that could be included in every robot using ROS and a kinect or other RGB-D sensor. This way, having different modularized functionalities it could be possible to create a customized robot in a few seconds, just importing packages and compiling them. 
\\[1cm]

\chapter{Context and motivation}
In the past fifteen years, there has been a change in the perception of the knowledge and whether it should be restricted using patents or not. Most of the scientific community now supports the open source initiative and this has caused a huge increase on the investigation field. \\



The idea of creating common software, of aiding other investigators to easily replicate the work already done and continue from that point instead of "reinventing the wheel" started in the past century. In 1998 the Open Source Initiative (OSI)\cite{osi} was formed, and its definition is recognized as the standard. According to them, "Open source does not just mean access to the source code. The distribution terms of open-source software must comply with the following criteria: Free Redistribution, Source Code, Derived Works, Integrity of The Author's Source Code, No Discrimination Against Persons or Groups, No Discrimination Against Fields of Endeavor, Distribution of License, License Must Not Be Specific to a Product, License Must Not Restrict Other Software,  License Must Be Technology-Neutral"\cite{osi_def}. 
\\

It is my believe that Open Source is critical in the development of new knowledge and not only in the Software field, but in all the science and technical disciplines. This is how the idea of increasing that common pool of tools with another one that might be useful appeared. 
\\
%\begin{wrapfigure}{r}{0.3\textwidth}
%	\centering
%   \includegraphics[width=0.25\textwidth]{img/intro/open_source.eps}
%	\caption[Open Source Initiative Logo]{Open Source Initiative Logo}
%\end{wrapfigure}

\begin{figure}[h]
	\begin{center}
    \includegraphics[scale=0.2]{img/intro/open_source.eps}
	\caption[Open Source Initiative Logo]{Open Source Initiative Logo}
	\end{center}
\end{figure}


It is noticeable the mark that Open Source has throughout the project. From the libraries being used in it to the Robotic Operating System of which this project is but another package. All of them are Open Source. Also, the tutorials, examples and different web-pages with useful comments and aids for those who are learning has been critical in this project's development. Nothing of this could have been possible without the idea of Open Source code. 
\\

%Apart from the want of making useful Open Source code, there remains another motivation unexplained. That is, the specific subject of this project. 
\\

%The bad economic context that exist nowadays has forced the different areas of the society to investigate new technologies that has a lower cost. 
%Nowadays, advances regarding robotics appear everyday, and technologies such as 3D printing are having a tremendous impact in the society. Thanks to this instruments, now it is possible and relatively easy to customize prosthesis and also much affordable than it was before. 
%\\

%Robotics is a field that is experiencing a huge impulse currently. And more specifically, Open Source robotics is growing rapidly. 
%\\

One of the open-source projects that appeared with this new way of thinking was the low-cost 3D printers. These machines have changed the manner of investigating in robots, since they allow to design them easily and have a 3D reproduction in a few hours. 
\\

But they are used not only in robots, but also in medicine and many other fields. In medicine, they have been a revolution since they allow to create customized pieces in very few time. They have been used for prosthesis and implants for persons of various ages, even for babies. 
\\

Before, the prosthesis were very expensive and permitted only fixed movements and combinations. The adaptations to each individual were made in the final product itself, trying to make it as comfortable as possible for the wearer. Nowadays, the prosthesis are customized for each patient, reducing the bothers for them and increasing their usability. Also, they can be easily and cheaply adapted for children as an example, who are still experimenting many changes in their bodies. They are much cheaper than they were before, and everyone with a 3D printer may construct one. 
\\

Of course, in order to 3D print a piece a file with its description is needed. There are many web-pages that store open-source designs that ranges from decoration models to complex prosthesis. This fact is decisive because there is not needed a huge amount of knowledge or money to improve the life quality of a person using these technologies. 
\\

There are numerous open-source projects and developers that put in common their knowledge to improve the technology being used. I have used many of them in the previous years, to learn about 3D printing, robotics or programming among other fields. 
\\

The fact that previously I have developed other projects involving designing and constructing 3D-printed wheeled robots gave me the idea of developing a new package for different robots that allowed to easily obtain new functionalities.One of the most interesting fields that can be used in virtually all types of robots is Computer Vision. 
\\

Computer vision has experience an important improvement in the last years through the upgrade of the hardware and software that compose it. The hardware such as acquisition elements (cameras, depth sensors, etc) and computing elements (PCs or other programmable devices) have experienced a rapid advance in the past years. It allowed to process more data that is now obtained more accurately and with less noise. This increase in the computing power of the equipment created a possibility of introducing more complex libraries and frameworks and even operating systems. 
\\

But Computer Vision is an ample area, why choosing object recognition? The project being developed is intended to be useful, not just a mere hobby. It occurred to me that apart from being useful in the interaction with robots and cognitive robots more particularly, this software could be used with visually impaired people, as an example. Mainly, with people that had a sudden loss of vision and that is not used to detect the objects by its shape, or cannot differentiate them by the contour. 
\\

In summary, the software developed is open-source and is intended as a useful package for visually impaired persons, as well as for its integration within a robotic system. 

\newpage
%\addcontentsline{toc}{part}{Introduction}
\part{Introduction}


%	
\section{Motivation}

Technology has evolved enormously in the past years. 
% In particular, robotics has changed and moved from controlled spaces such as factories to human inhabited spaces. 
% In section \ref{context} the reasons behind this shift in the robot's location and function were presented. 
% It was also estated that the importance of the perception systems augmented with the inclusion of robots in human-inhabited areas. 
% % This fact increases the importance of the perception systems being integrated. 
% The correct recognition of objects, persons, areas and situations is key for assistive and social tasks. 
Nowadays most of the robots being developed are only able to recognize a small part of their environment. 
% Most computer vision systems are currently used to navigate between points or to locate certain objects and grasp them. 
The decision algorithm of the robots is normally based on the instructions received from the user. 
Those commands are usually given by voice or inputting the desired task on a certain User Interface (UI).
\\

% For the robots to act as aiding personnel to take care of elder or sick people, the evolution of this decision algorithm is crucial. 
The algorithm decision of the robots must be upgraded if they are to perform tasks in a human inhabited area. %act as aiding personnel to take care of elder or sick people. 
They must be able to respond to commands, but also be aware of their environment and respond autonomously to it. 
% As an example, robots must be able to recognize the danger involved in certain objects or situations. 
% Also, they should perceive the environment and the interactions between humans to discern if, for example, they are arguing or if the conversation is not a beneficial one for the patient. 
% This latter example may be seen in the case of an anorexic person who is talking about food and weight with another person. 
% Then, the robot may intervene changing the subject subtly or trying to end the conversation. 
When this change has occurred, the robots may successfully perform the assistive tasks now reserved only to humans. 
But there is still a long way to research, mainly around the perception of the environment. 
The investment needed for its development is now being held due to the recession that appeared in the late 2000s decade, whose effects are nowadays still present. 
Nevertheless, this lack of funding might be mitigated using research, open data and open-source code. 
\\

I strongly believe in the ideals presented by the Open Source Initiative. 
This project is intended to be an open source code that can be a building block for other researchers that work on robot perception.
There are various scientists that have already studied the importance of the objects around the human in the action recognition. 
For example, in \cite{Delaitre}, a study of the human-object interaction in still images was performed in order to relate those objects to actions. 
	This relation between object and action may also be seen in \cite{Fathi}, in which wearable cameras are used to retrieve the input. 
	The objects being hand-held are recognized and the action associated to them is learned. 
	In this thesis I present a system that allows to easily learn and recognize hand-held objects in real time. 
	It is intended to be the previous step to the association between the object and the action. 



% The economic situation described above forced many experienced professionals to trying their luck creating start-ups. Many of the ideas of those enterprises are having nowadays a huge impact on the society. 
% \\

% One of these open-source projects that appeared was the low-cost 3D printers. These machines have changed the manner of investigating many fields, since they allow to design different pieces easily and have a 3D reproduction in a few hours. 
% \\

% Initially, they were used in investigation and more specifically in robotics, but now they are used in many different fields. Among those fields, there are medicine, construction or even food making. 
% \\

% In medicine, they have been a revolution since they allow to create customized and precise pieces in very few time. They have been used for prosthesis and implants for persons of various ages, even for babies. In the prosthesis fields in particular, the 3D printing technology is being a complete revolution. Before, the prosthesis were very expensive and permitted only fixed movements and combinations. The adaptations to each individual were made in the final product itself, trying to make it as comfortable as possible for the wearer. Nowadays, the prosthesis are customized for each patient, reducing the inconveniences and increasing their usability. Also, they can be easily and cheaply adapted for children as an example, who are still experimenting many changes in their bodies. They are much cheaper than they were before, and everyone with a 3D printer may construct one. 
% In order to 3D print a piece a file with its description is needed. There are many web-pages that store open-source designs that ranges from decoration models to complex prosthesis. This fact is decisive because there is not needed a huge amount of knowledge or money to improve the life quality of a person using these technologies. 
% \\

% There are numerous open-source projects and developers that put in common their knowledge to improve the technology being used. I have used many of them in the previous years, to learn about 3D printing, robotics or programming among other fields. 
% \\

% It is a fact that acquiring knowledge would be much difficult if the Open Source initiative has not been invented. This impulsed me towards developing something useful and that could be used by other people. The idea of creating a software that could be used in robotics investigation but also help people at the same time. 
% \\[0.5cm]

% Many of the projects are aimed at aiding physically impaired people, creating sternal skeletons and robotic arms that could aid them. But a personal fact led me to realize that visually impaired people were not having as much attention. The applications developed for them are still rough to use and also it is difficult for a grown person to develop his remaining senses to supply the information lost. 
% \\

% Besides, in the robotic field new lines of investigation have appeared. The social robots are now a reality and in the near future we will interact with them everyday. In order to understand the human behavior, the recognition of the objects being handled by them is crucial. 
% \\[0.5cm]


% Computer vision has experience an important improvement in the last years through the upgrade of the hardware and software that compose it. The hardware such as acquisition elements (cameras, depth sensors, etc) and computing elements (PCs or other programmable devices) have experienced a rapid advance in the past years. It allowed to process more data that is now obtained more accurately and with less noise. This increase in the computing power of the equipment created a possibility of introducing more complex libraries and frameworks and even operating systems. 
% \\

% Now, the technology is available to solve the problems presented, the aid of visually impaired people and the introduction of new information in the social robotics field. This is how the idea behind this thesis appeared: the creation of a modular software that implements an in-hand object recognition algorithm. 
 
%	%\addcontentsline{toc}{chapter}{Socio-economical context}
\section{Socio-economical context}

%Socio-economic - relating to both social and economic factors (social groups and the class system for example)
%Context - The circumstances/environment/events surrounding a specific thing. 



%	%\addcontentsline{toc}{chapter}{Future applications}
\chapter{Future applications}
	% future applications



\chapter{Introduction}
This bachelor's thesis consists on a software that implements an in-hand object learning and recognition using 2D and 3D information. 
\\

The code is Open-Source and it was designed to be running inside a robot. The idea is to create a new stand-alone module that could be included in every robot using ROS and a kinect or other RGB-D sensor. This way, having different modularized functionalities it could be possible to create a customized robot in a few seconds, just importing packages and compiling them. 
\\[1cm]

\chapter{Context and motivation}
In the past fifteen years, there has been a change in the perception of the knowledge and whether it should be restricted using patents or not. Most of the scientific community now supports the open source initiative and this has caused a huge increase on the investigation field. \\



The idea of creating common software, of aiding other investigators to easily replicate the work already done and continue from that point instead of "reinventing the wheel" started in the past century. In 1998 the Open Source Initiative (OSI)\cite{osi} was formed, and its definition is recognized as the standard. According to them, "Open source does not just mean access to the source code. The distribution terms of open-source software must comply with the following criteria: Free Redistribution, Source Code, Derived Works, Integrity of The Author's Source Code, No Discrimination Against Persons or Groups, No Discrimination Against Fields of Endeavor, Distribution of License, License Must Not Be Specific to a Product, License Must Not Restrict Other Software,  License Must Be Technology-Neutral"\cite{osi_def}. 
\\

It is my believe that Open Source is critical in the development of new knowledge and not only in the Software field, but in all the science and technical disciplines. This is how the idea of increasing that common pool of tools with another one that might be useful appeared. 
\\
%\begin{wrapfigure}{r}{0.3\textwidth}
%	\centering
%   \includegraphics[width=0.25\textwidth]{img/intro/open_source.eps}
%	\caption[Open Source Initiative Logo]{Open Source Initiative Logo}
%\end{wrapfigure}

\begin{figure}[h]
	\begin{center}
    \includegraphics[scale=0.2]{img/intro/open_source.eps}
	\caption[Open Source Initiative Logo]{Open Source Initiative Logo}
	\end{center}
\end{figure}


It is noticeable the mark that Open Source has throughout the project. From the libraries being used in it to the Robotic Operating System of which this project is but another package. All of them are Open Source. Also, the tutorials, examples and different web-pages with useful comments and aids for those who are learning has been critical in this project's development. Nothing of this could have been possible without the idea of Open Source code. 
\\

%Apart from the want of making useful Open Source code, there remains another motivation unexplained. That is, the specific subject of this project. 
\\

%The bad economic context that exist nowadays has forced the different areas of the society to investigate new technologies that has a lower cost. 
%Nowadays, advances regarding robotics appear everyday, and technologies such as 3D printing are having a tremendous impact in the society. Thanks to this instruments, now it is possible and relatively easy to customize prosthesis and also much affordable than it was before. 
%\\

%Robotics is a field that is experiencing a huge impulse currently. And more specifically, Open Source robotics is growing rapidly. 
%\\

One of the open-source projects that appeared with this new way of thinking was the low-cost 3D printers. These machines have changed the manner of investigating in robots, since they allow to design them easily and have a 3D reproduction in a few hours. 
\\

But they are used not only in robots, but also in medicine and many other fields. In medicine, they have been a revolution since they allow to create customized pieces in very few time. They have been used for prosthesis and implants for persons of various ages, even for babies. 
\\

Before, the prosthesis were very expensive and permitted only fixed movements and combinations. The adaptations to each individual were made in the final product itself, trying to make it as comfortable as possible for the wearer. Nowadays, the prosthesis are customized for each patient, reducing the bothers for them and increasing their usability. Also, they can be easily and cheaply adapted for children as an example, who are still experimenting many changes in their bodies. They are much cheaper than they were before, and everyone with a 3D printer may construct one. 
\\

Of course, in order to 3D print a piece a file with its description is needed. There are many web-pages that store open-source designs that ranges from decoration models to complex prosthesis. This fact is decisive because there is not needed a huge amount of knowledge or money to improve the life quality of a person using these technologies. 
\\

There are numerous open-source projects and developers that put in common their knowledge to improve the technology being used. I have used many of them in the previous years, to learn about 3D printing, robotics or programming among other fields. 
\\

The fact that previously I have developed other projects involving designing and constructing 3D-printed wheeled robots gave me the idea of developing a new package for different robots that allowed to easily obtain new functionalities.One of the most interesting fields that can be used in virtually all types of robots is Computer Vision. 
\\

Computer vision has experience an important improvement in the last years through the upgrade of the hardware and software that compose it. The hardware such as acquisition elements (cameras, depth sensors, etc) and computing elements (PCs or other programmable devices) have experienced a rapid advance in the past years. It allowed to process more data that is now obtained more accurately and with less noise. This increase in the computing power of the equipment created a possibility of introducing more complex libraries and frameworks and even operating systems. 
\\

But Computer Vision is an ample area, why choosing object recognition? The project being developed is intended to be useful, not just a mere hobby. It occurred to me that apart from being useful in the interaction with robots and cognitive robots more particularly, this software could be used with visually impaired people, as an example. Mainly, with people that had a sudden loss of vision and that is not used to detect the objects by its shape, or cannot differentiate them by the contour. 
\\

In summary, the software developed is open-source and is intended as a useful package for visually impaired persons, as well as for its integration within a robotic system. 

\newpage
%\addcontentsline{toc}{part}{Introduction}
\part{Introduction}


%	\input{intro/motivation} 
%	\input{intro/context}
%	\input{intro/uses}	% future applications



\chapter{Introduction}
This bachelor's thesis consists on a software that implements an in-hand object learning and recognition using 2D and 3D information. 
\\

The code is Open-Source and it was designed to be running inside a robot. The idea is to create a new stand-alone module that could be included in every robot using ROS and a kinect or other RGB-D sensor. This way, having different modularized functionalities it could be possible to create a customized robot in a few seconds, just importing packages and compiling them. 
\\[1cm]

\chapter{Context and motivation}
In the past fifteen years, there has been a change in the perception of the knowledge and whether it should be restricted using patents or not. Most of the scientific community now supports the open source initiative and this has caused a huge increase on the investigation field. \\



The idea of creating common software, of aiding other investigators to easily replicate the work already done and continue from that point instead of "reinventing the wheel" started in the past century. In 1998 the Open Source Initiative (OSI)\cite{osi} was formed, and its definition is recognized as the standard. According to them, "Open source does not just mean access to the source code. The distribution terms of open-source software must comply with the following criteria: Free Redistribution, Source Code, Derived Works, Integrity of The Author's Source Code, No Discrimination Against Persons or Groups, No Discrimination Against Fields of Endeavor, Distribution of License, License Must Not Be Specific to a Product, License Must Not Restrict Other Software,  License Must Be Technology-Neutral"\cite{osi_def}. 
\\

It is my believe that Open Source is critical in the development of new knowledge and not only in the Software field, but in all the science and technical disciplines. This is how the idea of increasing that common pool of tools with another one that might be useful appeared. 
\\
%\begin{wrapfigure}{r}{0.3\textwidth}
%	\centering
%   \includegraphics[width=0.25\textwidth]{img/intro/open_source.eps}
%	\caption[Open Source Initiative Logo]{Open Source Initiative Logo}
%\end{wrapfigure}

\begin{figure}[h]
	\begin{center}
    \includegraphics[scale=0.2]{img/intro/open_source.eps}
	\caption[Open Source Initiative Logo]{Open Source Initiative Logo}
	\end{center}
\end{figure}


It is noticeable the mark that Open Source has throughout the project. From the libraries being used in it to the Robotic Operating System of which this project is but another package. All of them are Open Source. Also, the tutorials, examples and different web-pages with useful comments and aids for those who are learning has been critical in this project's development. Nothing of this could have been possible without the idea of Open Source code. 
\\

%Apart from the want of making useful Open Source code, there remains another motivation unexplained. That is, the specific subject of this project. 
\\

%The bad economic context that exist nowadays has forced the different areas of the society to investigate new technologies that has a lower cost. 
%Nowadays, advances regarding robotics appear everyday, and technologies such as 3D printing are having a tremendous impact in the society. Thanks to this instruments, now it is possible and relatively easy to customize prosthesis and also much affordable than it was before. 
%\\

%Robotics is a field that is experiencing a huge impulse currently. And more specifically, Open Source robotics is growing rapidly. 
%\\

One of the open-source projects that appeared with this new way of thinking was the low-cost 3D printers. These machines have changed the manner of investigating in robots, since they allow to design them easily and have a 3D reproduction in a few hours. 
\\

But they are used not only in robots, but also in medicine and many other fields. In medicine, they have been a revolution since they allow to create customized pieces in very few time. They have been used for prosthesis and implants for persons of various ages, even for babies. 
\\

Before, the prosthesis were very expensive and permitted only fixed movements and combinations. The adaptations to each individual were made in the final product itself, trying to make it as comfortable as possible for the wearer. Nowadays, the prosthesis are customized for each patient, reducing the bothers for them and increasing their usability. Also, they can be easily and cheaply adapted for children as an example, who are still experimenting many changes in their bodies. They are much cheaper than they were before, and everyone with a 3D printer may construct one. 
\\

Of course, in order to 3D print a piece a file with its description is needed. There are many web-pages that store open-source designs that ranges from decoration models to complex prosthesis. This fact is decisive because there is not needed a huge amount of knowledge or money to improve the life quality of a person using these technologies. 
\\

There are numerous open-source projects and developers that put in common their knowledge to improve the technology being used. I have used many of them in the previous years, to learn about 3D printing, robotics or programming among other fields. 
\\

The fact that previously I have developed other projects involving designing and constructing 3D-printed wheeled robots gave me the idea of developing a new package for different robots that allowed to easily obtain new functionalities.One of the most interesting fields that can be used in virtually all types of robots is Computer Vision. 
\\

Computer vision has experience an important improvement in the last years through the upgrade of the hardware and software that compose it. The hardware such as acquisition elements (cameras, depth sensors, etc) and computing elements (PCs or other programmable devices) have experienced a rapid advance in the past years. It allowed to process more data that is now obtained more accurately and with less noise. This increase in the computing power of the equipment created a possibility of introducing more complex libraries and frameworks and even operating systems. 
\\

But Computer Vision is an ample area, why choosing object recognition? The project being developed is intended to be useful, not just a mere hobby. It occurred to me that apart from being useful in the interaction with robots and cognitive robots more particularly, this software could be used with visually impaired people, as an example. Mainly, with people that had a sudden loss of vision and that is not used to detect the objects by its shape, or cannot differentiate them by the contour. 
\\

In summary, the software developed is open-source and is intended as a useful package for visually impaired persons, as well as for its integration within a robotic system. 

\newpage
\input{intro/intro} 
\newpage
\input{intro/thesis_structure}
 
\newpage
\addcontentsline{toc}{chapter}{Thesis structure}
\chapter*{ Thesis structure}

 
\newpage
\addcontentsline{toc}{chapter}{Thesis structure}
\chapter*{ Thesis structure}



%%%%%% PART 1 %%%%%%

%section or whatever

OpenCV\cite{opencv} (Open Source Computer Vision) is a library that implements real-time computer vision algorithms. 

It is cross-platform and it is released under a BSD\cite{BSD} license.

\begin{figure}[h]
	\begin{center}
    \includegraphics[scale=1]{img/opencv/logo.eps}
	\caption[OpenCV Logo]{OpenCV Logo}
	\end{center}
\end{figure}


\include {pcl}

"The Robot Operating System (ROS) is a set of software libraries and tools that help you build robot applications. From drivers to state-of-the-art algorithms, and with powerful developer tools, ROS has what you need for your next robotics project. And it's all open source."\cite{ros}

\begin{figure}[h]
	\begin{center}
    \includegraphics[scale=0.3]{img/ros/groovy.eps}
	\caption[ROS Groovy Logo]{ROS Groovy Logo}
	\end{center}
\end{figure}

%%%%%% BIBLIOGRAPHY %%%%%%
\newpage
\bibliography{tfg}{}	%poner el archivo .bib
\bibliographystyle{plain}


\end{document}
